\Artigo%
{divulgacion}%
{Sobre a lóxica da física}%
{Ana Peón Nieto}%
{}%

\begin{multicols}{2}

A ninguén sorprenderá se digo que este é un artigo sobre as matemáticas da
física. Algúns pasarán a páxina con preguiza, outros esperarán ecuacións
formais moi bonitas e globais, imposibles de manipular (isto vai polos meus
pobres estudantes de Métodos II,  compréndovos!). Pois non. En absoluto. Este
artigo fala de lóxica matemática, unha das ramas máis teóricas, descoñecidas, e
(desgraciadamente) ignoradas polos demais profesionais das matemáticas (salvo
excepcións, son coma as bruxas, habelas hainas). É, por así dicilo, a teoría de
cordas das mates.

O curioso do asunto é que hai moi boas cabezas pensando nas aplicacións desta
disciplina á física. Non tanto como ferramenta formal, senón para resolver
cuestións fundamentais. Un pouco como no caso da teoría de cordas, que xurdiu
coa motivación de unificar a gravidade e a mecánica cuántica; a lóxica
matemática, e máis precisamente a subdisciplina chamada teoría de modelos,
inquírese sobre a base das teorías: cal é a linguaxe adecuada? Cales son os
axiomas (léase postulados) que garanten o sustento da teoría? A moitos soará o
teorema de incompletitude de Gödel, segundo o cal, en «matemáticas», entendida
coma a ciencia na que os números enteiros son base, sempre haberá enunciados
indecidibles (é dicir, cuxa veracidade ou falsidade non se poida demostrar).

Pero: que é a teoría de modelos? Que a fai adecuada para esta inmensa tarefa?
Pois ben, dunha parte, a teoría de modelos ocúpase de calquera cousa
expresable matematicamente. Por exemplo, os grupos abelianos, ou conmutativos,
coma $\mathbb{Z}$. Para isto, escollemos unha linguaxe que inclúa a función +,
os cuantificadores $\forall,\, \exists$, e un signo especial para o cero. Con
isto podemos establecer os axiomas ou normas básicas:
%
\begin{align*}
    &\forall x,y          & & x+y=y+x,\\
    &\forall y            & & 0+y=y+0=y,\\
    &\forall x\ \exists y & & x+y=y+x=0,\\
    &\forall x,y,z        & & (x+y)+z=x+(y+z).
\end{align*}
%
Dende o punto de vista da teoría de modelos, un grupo abeliano é calquera
«universo» ou «estrutura» matemática que satisfaga estes axiomas. Dá igual
que teña un elemento, cinco ou infinitos. Estas simples normas capturan o
esencial. (Non abandonedes aínda, por favor!). O máis interesante é que
considerando estes axiomas, o lóxico ten tódolos posibles grupos na cabeza ao
mesmo tempo. E o mesmo aplica a calquera teoría das matemáticas: para o lóxico
non é cuestión deste ou dese espazo, senón de todos á vez. Isto, que pode
parecer pouco práctico, dá lugar a aplicacións impresionantes en análise, por
exemplo, ou xeometría. E ademais poderedes contestar aos profesores de Métodos
cando vos critiquen por riscardes o diferencial $\text{d}x$. En efecto, na teoría dos
corpos ordenados, coma os números reais, existe un universo no que o
diferencial é un número coma outro calquera, máis grande ca $0$, pero máis
pequeno que tódolos reais positivos. Así que a riscar diferenciais sen
complexo!

Pois ben, perdidos andaban os lóxicos nas súas disquisicións existenciais,
cando Hrushovski e Zilber, dous lóxicos brillantes\footnote{O nome do primeiro
barallouse para a Medalla Fields, o máis alto galardón en matemáticas (que, por
certo, ostenta Witten como único físico na lista). Lamentablemente, na miña
humilde opinión, Hrushovski non levou este premio. De llo ter dado, teríase
visibilizado o seu enorme labor e achegado a lóxica ao matemático medio.},
deron coas famosas estruturas de Zariski. Buscaban \textbf{A} xeometría.
\textbf{A xeometría}. E deron con algo estraño aos ollos dos comúns, é dicir,
uns modelos que capturaban á vez a xeometría do espazo macroscópico, no que
mirar nunha orde ou noutra non afecta a observación (é dicir, é conmutativa),
pero que eran non conmutativos en esencia. Coma o mundo cuántico. Disto hai
moitos anos, a finais do século XX xa se aplicaban estes modelos para probar
teoremas importantes, coma a conxectura de Mordell-Lang \cite{Ehud_1996}.

Pero Zilber, apaixonado da física, que perseguía una formulación lóxica (en
tódolos sentidos) desta ciencia, viuno claro. E se podemos aproximar o mundo
real, incluíndo o mundo macroscópico e cuántico, mediante estas estruturas? En
colaboración con outros dous lóxicos, Solanki e Sustretov, estudaron o caso
máis sinxelo: o oscilador harmónico unidimensional. Hoxe centrareime neste
traballo, pero deixádeme engadir que Zilber segue no seu empeño de atopar
\textbf{A} linguaxe (\textbf{A LINGUAXE}) da física dende a teoría de modelos.
Nos últimos anos, catro artigos prometedores axiomatizan a mecánica cuántica (e
a unifican coa cuántica estatística) mediante a lóxica do continuo
\cite{Zilber1_2025}\cite{Zilber2_2025}\cite{Zilber3_2024}\cite{Zilber_2023}.
Espero comprendelos nalgún momento para explicalo como hoxe farei co oscilador.

Comecemos polo principio: o oscilador harmónico en dimensión 1 (é dicir, un
resorte cun extremo ligado a un punto, o marco, e outro que se move libremente).
A enerxía total do sistema é
%
\begin{equation}
    E=\frac{m}{2}\dot{x}^2+\frac{k}{2}x^2,
\end{equation}
%
onde $V(x)=\frac{k}{2}x^2$ é a enerxía potencial, con $k$ a constante elástica. Polo principio de conservación da enerxía, empregando algún
truquiño de integración, obtense a ecuación
%
\begin{equation}
    m\ddot{x}+kx=0,
\end{equation}
%
cuxa solución xeral é $x(t)=\cos(\omega t+\alpha)$, con
$\omega=\sqrt{\frac{k}{m}}$. Podemos resolver no plano considerando pares
$(x(t),\dot{x}(t))$. Se chamamos $p(t)=m\dot{x}(t)$ (é dicir, o momento
lineal), entón temos
%
\begin{equation}
    \dot{x}(t)=\frac{p(t)}{m},\qquad\dot{p}(t)=\frac{k}{m}.
\end{equation}
%
Para cuantizar, imos poñer todo en notación hamiltoniana. Temos as coordenadas
canónicas $p(t)$ e $q(t)$ (que neste caso son tan só o momento lineal
$p(t)=m\dot{x}(t)$ e a posición $q(t)=x(t)$), nas que obtemos as ecuacións de
Hamilton (que son as mesmas que antes, pero son válidas nun contexto máis
amplo, para tódolos sistemas hamiltonianos)
%
\begin{equation}
    \dot{p}=\{H,p\}=-\frac{\partial H}{\partial q},
    \qquad
    \dot{q}=\{H,q\}=\frac{\partial H}{\partial p}
\end{equation}
%
onde $H(p,q)=\frac{p^2}{2m}+\frac{m\omega^2q^2}{2}$ é o Hamiltoniano, neste
caso, a enerxía total do sistema, e a diferenza de signos nas ecuacións
débese a xeometría simpléctica subxacente. Os corchetes $\{\cdot \}$ chámanse
corchetes de Poisson, e teñen sentido para tódalas funcións
$\{f(p,q),g(p,q)\}$; o que mide é a diferenza das variacións de $f$ e $g$ en 
certas direccións relacionadas co momento e posición.

Á hora de cuantizar, os observables, neste caso $p$ e $q$, pasan a
interpretarse como operadores actuando no espazo de Hilbert dos estados do
sistema, neste caso, o espazo vectorial $\mathbb{C}$, adornado cunha estrutura
hermítica. A avaliación de operadores interprétase como a medición. Se $P, Q$
son os observables cuantizados, escribimos $[P,Q]=PQ-QP$. Eses corchetes son a
cuantización dos corchetes de Poisson. Aplicar $PQ$ a un estado do sistema
corresponde a observar primeiro a posición e, inmediatamente despois, o
momento. Sabemos que no mundo cuántico, a observación modifica o estado do
sistema, máis explicitamente, $[P,Q]=-i\hbar$ onde $\hbar$ é a constante de
Planck. Temos o Hamiltoniano cuántico
%
\begin{equation}
    H=\frac{P^2}{2m}+\frac{m\omega^2Q^2}{2}
\end{equation}
%
que pertence á «álxebra» de observables (unha forma curta de dicir que
combinamos observacións mediante consecución e suma). Entón, as solucións
$\psi(x)$ da ecuación diferencial
%
\begin{equation}
    H\psi(x)=\hbar \omega\left(n+\frac{1}{2}\right) \psi(x)
\end{equation}
%
para $n\in\mathbb{N}$ comprenden tódolos posibles estados enerxéticos do
sistema. Así mesmo, podemos pasar dun estado $|n\rangle$ a un estado
$|n+1\rangle$ e viceversa mediante os operadores de creación
$a^\dagger=\sqrt{\frac{m\omega}{2k}}Q-\frac{i}{\sqrt{2\hbar\omega}}P$ e
aniquiliación $a=\sqrt{\frac{m\omega}{2k}}Q+\frac{i}{\sqrt{2\hbar\omega}}P$,
respectivamente.

Tras este recordatorio, volvamos ao noso, aínda que de maneira un tanto
camuflada. A idea dos tres lóxicos é lóxica: para cada $N\in\mathbb{N}$,
consideran un espazo $L_N$ consistente nunha recta $\mathbb{A}$ (quizais co
infinito $\mathbb{P}^1=\mathbb{A}\cup\{\infty\}$, pensade nos reais co
infinito, ou os complexos co infinito; en realidade, pensade en todos á vez,
porque uns son os puntos reais da recta, e os outros, os puntos complexos) e, sobre cada
punto $x\in\mathbb{A}$, o grupo de raíces $N$-ésimas da unidade (ou
$\mathbb{Z}_N$ se o preferides). Estes son os números $z$ tal que $z^N=1$. Por
exemplo, se $N=2$, $z=\pm1$, que observamos que corresponden aos valores
$e^{\frac{2k\pi i}{2}}$ para $k=0,1$. Para $N$ xeral, obteremos $e^{\frac{2k\pi
i}{N}}$ para $k=0,\dots, N-1$. Este grupo $\mathbb{Z}_N$ (é un grupo
esencialmente porque se poden multiplicar elementos e seguimos dentro, o $1$
pertence a $\mathbb{Z}_N$, e ademais tódolos elementos teñen inverso)
representa os posibles estadíos do sistema no punto $x$ sobre o que vive, ata
enerxía $N$. Para conseguir tódolos estadíos enerxéticos haberá que
considerar tódolos $N$'s xuntos\footnote{
Isto non é problemático, pois coma veredes unha vez acabemos, as construcións
son compatibles coa relación $N<N+1$, no sentido de que $L_N \subset L_{N+1}$,
e as operacións de creación e aniquilación de $L_{N+1}$ inducen as de $L_N$.}.
Obsérvese que temos
%
\begin{equation}
    \pi:L_N=\mathbb{P}^1\times\mathbb{Z}_N\longrightarrow \mathbb{P}^1
\end{equation}
%
que a cada punto $\ell=(x,\gamma)\in L_N$ asígnalle o $x\in \mathbb{P}^1$ sobre
o que «vive». Así, a fibra $\pi^{-1}(x)$ representaría os estados do sistema
en $x$.

Os operadores de creación (resp. aniquilación) pódense modelizar considerando
tripletes $(\ell,\ell',b)\in A^\dagger\subset L_N\times L_N\times \mathbb{A}$
con $\ell\in\pi^{-1}(x)$, $\ell'\in\pi^{-1}(x+1)$, $b^2=x$ (de maneira que
$(\gamma\ell',\gamma\ell,b_\gamma)\in A^\dagger$ para todo
$\gamma\in\mathbb{Z}_N$ e algún $b_\gamma\in\mathbb{A}$)  (resp.
$(\gamma\ell',\gamma\ell,b_\gamma)\in A\subset L_N\times L_N\times \mathbb{A}$
con $\ell'\in\pi^{-1}(x+1)$, $\ell\in\pi^{-1}(x)$, $b_\gamma^2=x$). Se o
triplete $(\ell,\ell',b)\in A^\dagger$, interpretamos que
pasamos do estado $\ell$ (que identificamos con algún $|M\rangle$) no punto $x$
ao estado $\ell'\sim|M+1\rangle$.
Coma en principio os estados represéntanse mediante puntos da recta, podemos
identificar o estado $x$ con $|M\rangle$, e $x+1$ con $|M+1\rangle$. O grupo vén
a conto porque en principio, o oscilador podería encontrarse noutro estado
enerxético no punto $x$, polo que $\gamma\ell$ representaría un sistema no
que $x$ é $|M+\frac{2\pi i}{N}\log\gamma\rangle$. A vantaxe desta linguaxe é
que non necesitamos identificar $M$ exactamente, só saber que o sistema está
nalgún estado, en consonancia coa cuántica\footnote{
A raíz $b^2=x$ ten que ver con como actúan os operadores cuánticos, é 
importante, e de feito é o \textit{quid} desta construción tamén.}.

Hai unha certa duplicidade na definición, xa que o feito de crear está na
aparición de $\ell$ e $\ell'$ no mesmo triplete (sobre puntos diferentes), pero
usamos o punto da base para levar a conta de cantos estados enerxéticos
modificamos. Como basta saber que estados se modifican nun paso, e despois
noutro, etc, basta considerar o estado inicial e o punto final. De feito, os 
autores dan outra construción posible (e equivalente nun sentido adecuado
\cite{Solanki_2013}).

Identificamos, dentro de $L_N\times\mathbb{A}$, os
puntos $(\ell\gamma,\gamma^{-1}y)$ para todo $\gamma\in\mathbb{Z}$; é dicir,
consideramos o conxunto
$[(\ell,y)]=\{(\ell\gamma,\gamma^{-1}y)\colon\gamma\in\mathbb{Z}\}$ como un
único punto. Entón, o conxunto $\mathcal{H}_x=\{[\ell,y]\colon \pi(\ell)=x\}$ é
un espazo vectorial para o produto por escalares $t\cdot [\ell,y]=[\ell,ty]$ e
suma
%
\begin{align}
    [\ell_1,y_1]+[\ell_2,y_2]\stackrel{\tiny\exists[\ell_1,y_2']=[\ell_2,y_2]}{=}
    &[\ell_1,y_1]+[\ell_1,y_2']\\
    =&[\ell_1,y_1+y_2],
    \footnote{Exercicio: comproba que está ben definido!}
\end{align}
%
sobre o que se poden definir os operadores
%
\begin{equation}
0    a^\dagger[(e,y)]=[(e',yz)],\qquad a[(e',y)]=[(e,yz)]
\end{equation}
%
onde $(e,e',z)\in A^\dagger$, ou $(e',e,y)\in A$. Ben, a cuestión é que
considerar $A^\dagger$ ou $a^\dagger$, $A$ ou $a$, non cambia a lóxica do
asunto, pois as teorías de cada un destes espazos (é dicir, tódolos teoremas
satisfeitos por cada un deles) son equivalentes nun sentido adecuado que
permite obter unha da outra.

E agora me diredes: pero isto non é xeometría? Onde está a lóxica? Pois a
lóxica está en que nada disto se pode enunciar formalmente no mundo da
xeometría \textbf{complexa} clásica. Pódese se utilizamos os reais (a clave é o
uso dunha raíz cadrada na definición das relacións, cuxas dúas solucións
reais, positiva ou negativa, son identificables coa creación ou aniquilación).
Pode que isto non pareza revolucionario, pero espérase que os exemplos en
teoría de modelos veñan da xeometría complexa, e este non o cumpre. O que si
que é máis interesante é que $L_N$ vive sobre a recta, que \textbf{SI} que é
complexa, no sentido de que podemos coller puntos reais ou complexos sen
cambiarmos nada. Un recubrimento finito dunha recta que non é complexo (coma o
$L_N$ aquí presente) é \textbf{MOI} insólito. A súa existencia débese á
subxacencia de xeometría non conmutativa, que se percibe no feito de que os 
operadores de creación e aniquilación non conmutan, do mesmo modo que a posición 
e o momento cuánticos non conmutan. De feito, se
miramos a obxectos máis refinados, coma os automorfismos do recubrimento
$L_N\longrightarrow \mathbb{A}$, veremos que o comportamento é estraño, e
codifica a relación $[P,Q]=-i\hbar$. Para construírmos espazos desta índole,
temos que movernos ao mundo das xeometrías de Zariski puras, fundamentalmente
modelo-teóricas, non provenientes de ningún espazo clásico.

Recapitulando: a lóxica permítenos construír un espazo, que pode parecer
xeométrico, pero non o é, onde vive o oscilador harmónico cuántico. Aínda que é
un exemplo sinxelo, é tamén o primeiro exemplo físico que se consegue modelizar
completamente con ferramentas lóxicas. Así mesmo, prové un dos primeiros
exemplos dunha xeometría de Zariski non proveniente da xeometría clásica, o cal
engade interese a estas teorías. Non sei a vós, pero a min éncheme de
expectación...

\printbibliography
\end{multicols}
